\chapter{Metode deret pangkat} \label{ps:chapter}

%%%%%%%%%%%%%%%%%%%%%%%%%%%%%%%%%%%%%%%%%%%%%%%%%%%%%%%%%%%%%%%%%%%%%%%%%%%%%%

\section{Deret pangkat}
\label{powerseries:section}

%mbxINTROSUBSECTION

\sectionnotes{1,5 atau 2 kuliah\EPref{, \S8.1 dalam \cite{EP}}\BDref{,
\S5.1 dalam \cite{BD}}}

Banyak fungsi dapat dituliskan dalam bentuk deret pangkat
\begin{equation*}
\sum_{k=0}^\infty a_k {(x-x_0)}^k .
\end{equation*}
Dengan mengandaikan bahwa solusi suatu persamaan diferensial merupakan
deret pangkat, kita mungkin dapat menggunakan metode yang mengingatkan pada koefisien
tak tentu---kita mencoba menentukan bilangan-bilangan $a_k$.
Sebelum menjalankan proses ini, kita meninjau beberapa hasil
dan konsep tentang deret pangkat.

\subsection{Definisi}

Seperti telah kita katakan, suatu \emph{\myindex{deret pangkat}} adalah ekspresi seperti
\begin{equation} \label{ps:sereq1}
\sum_{k=0}^\infty a_k {(x-x_0)}^k =
a_0 + 
a_1 (x-x_0) +
a_2 {(x-x_0)}^2 +
a_3 {(x-x_0)}^3 + \cdots,
\end{equation}
di mana $a_0,a_1,a_2,\ldots,a_k,\ldots$ dan $x_0$ adalah konstanta. Misalkan
\begin{equation*}
%mbxlatex \begin{split}
S_n(x)
%mbxlatex &
=
\sum_{k=0}^n a_k {(x-x_0)}^k
%mbxlatex \\
%mbxlatex &
=
a_0 + a_1 (x-x_0) + a_2 {(x-x_0)}^2 + a_3 {(x-x_0)}^3 + \cdots + a_n {(x-x_0)}^n ,
%mbxlatex \end{split}
\end{equation*}
menyatakan apa yang disebut \emph{\myindex{jumlah parsial}}. Jika untuk suatu $x$,
limit
\begin{equation*}
\lim_{n\to \infty} S_n(x) = \lim_{n\to\infty} \sum_{k=0}^n a_k {(x-x_0)}^k
\end{equation*}
ada, kita mengatakan bahwa deret \eqref{ps:sereq1}
\emph{konvergen}\index{konvergensi deret pangkat} di $x$.
Di $x=x_0$, deret selalu konvergen ke $a_0$.
Jika \eqref{ps:sereq1}
konvergen di sembarang $x \neq x_0$ lainnya,
kita mengatakan bahwa \eqref{ps:sereq1} adalah suatu
\emph{\myindex{deret pangkat konvergen}}, dan kita menulis
\begin{equation*}
\sum_{k=0}^\infty a_k {(x-x_0)}^k = 
\lim_{n\to\infty} \sum_{k=0}^n a_k {(x-x_0)}^k.
\end{equation*}
Jika deret tersebut tidak konvergen di titik mana pun $x \neq x_0$, kita mengatakan bahwa
deret tersebut \emph{divergen}\index{deret pangkat divergen}.

\begin{example} \label{ps:expex}
Deret
\begin{equation*}
\sum_{k=0}^\infty \frac{1}{k!} x^k = 
1 + x + \frac{x^2}{2} + \frac{x^3}{6} + \cdots
\end{equation*}
konvergen untuk setiap $x$.
Ingat bahwa $k! = 1\cdot 2\cdot 3 \cdots k$ adalah
faktorial. Berdasarkan konvensi, kita mendefinisikan $0! = 1$.
Anda mungkin ingat bahwa deret ini
konvergen ke $e^x$.
\end{example}

Kita mengatakan bahwa \eqref{ps:sereq1}
\emph{\myindex{konvergen mutlak}}\index{konvergensi mutlak}
di $x$ apabila limit
\begin{equation*}
\lim_{n\to\infty} \sum_{k=0}^n
\lvert a_k \rvert \, {\lvert x-x_0 \rvert}^k 
\end{equation*}
ada. Artinya, deret
$\sum_{k=0}^\infty \lvert a_k \rvert \, {\lvert x-x_0 \rvert}^k$
konvergen.
Jika \eqref{ps:sereq1} konvergen mutlak di $x$, maka deret tersebut
konvergen di $x$. Namun, implikasi sebaliknya tidak benar.

\begin{example} \label{ps:1kex}
Deret
\begin{equation*}
\sum_{k=1}^\infty \frac{1}{k} x^k
\end{equation*}
konvergen mutlak untuk semua $x$ dalam interval $(-1,1)$.
Deret tersebut konvergen di $x=-1$,
karena
$\sum_{k=1}^\infty \frac{{(-1)}^k}{k}$ konvergen (secara bersyarat)
berdasarkan uji deret
berselang-seling.
Deret pangkat tersebut tidak konvergen mutlak di $x=-1$, karena
$\sum_{k=1}^\infty \frac{1}{k}$ tidak konvergen.
Deret tersebut
divergen di $x=1$.
\end{example}

\subsection{Jari-jari konvergensi}

Jika suatu deret pangkat konvergen mutlak
di suatu $x_1$, maka untuk semua $x$ sedemikian sehingga
$\lvert x - x_0  \rvert \leq \lvert x_1 - x_0 \rvert$ (yaitu, $x$ tidak
lebih jauh dari $x_0$ daripada $x_1$), kita mempunyai
$\lvert a_k {(x-x_0)}^k \rvert \leq
\lvert a_k {(x_1-x_0)}^k \rvert$
untuk semua $k$.
Karena bilangan-bilangan $\lvert a_k {(x_1-x_0)}^k \rvert$ berjumlah hingga suatu
limit berhingga, penjumlahan bilangan-bilangan positif yang lebih kecil
$\lvert a_k {(x-x_0)}^k \rvert$ juga harus mempunyai limit berhingga.
Oleh karena itu, deret tersebut harus konvergen
mutlak di $x$. %  Kita mempunyai hasil berikut.

\begin{theorem}
Untuk suatu deret pangkat \eqref{ps:sereq1}, terdapat bilangan
$\rho$ (kita mengizinkan $\rho=\infty$)
yang disebut \emph{\myindex{jari-jari konvergensi}} sedemikian sehingga
deret tersebut konvergen mutlak pada interval
$(x_0-\rho,x_0+\rho)$ dan divergen untuk $x < x_0-\rho$ dan $x > x_0+\rho$.
Kita menulis $\rho=\infty$ jika
deret tersebut konvergen untuk semua $x$.
\end{theorem}

\begin{myfig}
\capstart
\inputpdft{ps-conv}{%
Diagram garis bilangan dengan titik-titik x sub 0 dikurangi rho,
x sub 0, dan x sub 0 ditambah rho ditandai.
Ruas antara x sub 0 dikurangi rho dan x sub 0 ditambah rho diarsir dan ditandai
'konvergen mutlak'. Bagian garis di bawah x sub 0 dikurangi rho dan
di atas x sub 0 ditambah rho diberi label 'divergen'.}
\caption{Konvergensi suatu deret pangkat.\label{ps:convfig}}
\end{myfig}

Lihat \figurevref{ps:convfig}.
Dalam \exampleref{ps:expex}, jari-jari konvergensinya adalah $\rho = \infty$
karena deret tersebut konvergen di mana-mana.
Dalam \exampleref{ps:1kex}, jari-jari konvergensinya adalah $\rho=1$.
Kita mencatat bahwa $\rho = 0$ adalah cara lain untuk mengatakan bahwa deret tersebut
divergen.

Uji yang berguna untuk konvergensi suatu deret adalah
\emph{uji rasio}\index{uji rasio untuk deret}. Andaikan bahwa
\begin{equation*}
\sum_{k=0}^\infty c_k
\end{equation*}
adalah suatu deret dan limit
\begin{equation*}
L = \lim_{k\to\infty} \left \lvert \frac{c_{k+1}}{c_k} \right \rvert
\end{equation*}
ada. Maka deret tersebut konvergen mutlak jika $L < 1$ dan divergen
jika $L > 1$.

Kita menerapkan uji ini pada deret pangkat \eqref{ps:sereq1}. 
Tetapkan $c_k = a_k {(x - x_0)}^k$ dalam uji tersebut. Hitung
\begin{equation*}
L = \lim_{k\to\infty} \left \lvert \frac{c_{k+1}}{c_k} \right \rvert
=
\lim_{k\to\infty} \left \lvert
\frac{a_{k+1} {(x - x_0)}^{k+1}}{a_k {(x - x_0)}^k}
\right \rvert
=
\lim_{k\to\infty} \left \lvert
\frac{a_{k+1}}{a_k}
\right \rvert
\lvert  x - x_0 \rvert .
\end{equation*}
Definisikan $A$ dengan
\begin{equation*}
A =
\lim_{k\to\infty} \left \lvert
\frac{a_{k+1}}{a_k}
\right \rvert .
\end{equation*}
Maka deret \eqref{ps:sereq1}
konvergen mutlak jika $1 > L = A \lvert x - x_0 \rvert$.
Jika $A > 0$, maka
deret tersebut konvergen mutlak
jika $\lvert x - x_0 \rvert < \nicefrac{1}{A}$,
dan divergen jika $\lvert x - x_0 \rvert > \nicefrac{1}{A}$. Artinya,
jari-jari konvergensinya adalah $\nicefrac{1}{A}$.
Jika $A = 0$, maka deret tersebut selalu konvergen.

\emph{Uji akar} serupa\index{uji akar untuk deret}.
Andaikan
\begin{equation*}
L = \lim_{k\to\infty} \sqrt[k]{\lvert c_k \rvert}
\end{equation*}
ada. Maka $\sum_{k=0}^\infty c_k$ konvergen mutlak jika $L < 1$
dan divergen jika $L > 1$. Kita dapat menggunakan perhitungan yang sama seperti di atas
untuk mencari $A$.
Mari kita rangkum.

\begin{theorem}[Uji rasio dan uji akar untuk deret pangkat]
Tinjau suatu deret pangkat
\begin{equation*}
\sum_{k=0}^\infty a_k {(x-x_0)}^k
\end{equation*}
sedemikian sehingga
\begin{equation*}
A =
\lim_{k\to\infty}
\left \lvert
\frac{a_{k+1}}{a_k}
\right \rvert
\qquad \text{atau} \qquad
A =
\lim_{k\to\infty} \sqrt[k]{\lvert a_k \rvert}
\end{equation*}
ada. Jika $A = 0$, maka jari-jari konvergensi deret tersebut
adalah $\infty$. Jika tidak, jari-jari konvergensinya adalah $\nicefrac{1}{A}$.
Selain itu, jika salah satu limit untuk 
$A$ divergen ke $\infty$, maka deret tersebut divergen.
\pagebreak[3]
\end{theorem}

\begin{example}
Tinjau
\begin{equation*}
\sum_{k=0}^\infty 2^{-k} {(x-1)}^k .
\end{equation*}
Kita menghitung limit dalam uji rasio,
\begin{equation*}
A = \lim_{k\to\infty} 
\left \lvert
\frac{a_{k+1}}{a_k}
\right \rvert
=
\lim_{k\to\infty} 
\left \lvert
\frac{2^{-k-1}}{2^{-k}}
\right \rvert
=
\lim_{k\to\infty} 
2^{-1} = \frac{1}{2}.
\end{equation*}
Oleh karena itu, jari-jari konvergensinya adalah $2$, dan deret tersebut
konvergen mutlak pada interval $(-1,3)$.
Kita juga dapat menggunakan uji akar:
\begin{equation*}
A = 
\lim_{k\to\infty} 
\sqrt[k]{\lvert
a_k
\rvert}
=
\lim_{k\to\infty} 
\sqrt[k]{\lvert
2^{-k}
\rvert}
=
\lim_{k\to\infty} 
2^{-1}
=
\frac{1}{2}.
\end{equation*}
\end{example}

\begin{example}
Tinjau
\begin{equation*}
\sum_{k=1}^\infty \frac{1}{k^k} {x}^k .
\end{equation*}
Hitung limit untuk uji akar (uji rasio juga dapat digunakan),
\begin{equation*}
A =
\lim_{k\to\infty} 
\sqrt[k]{\lvert a_k \rvert}
=
\lim_{k\to\infty} 
\sqrt[k]{
\left\lvert\frac{1}{k^k}\right\rvert}
=
\lim_{k\to\infty} 
\sqrt[k]{
{\left\lvert\frac{1}{k}\right\rvert}^{k}}
=
\lim_{k\to\infty} 
\frac{1}{k}
=
0 .
\end{equation*}
Jadi jari-jari konvergensinya adalah $\infty$: Deret tersebut
konvergen di mana-mana. %  Uji rasio juga dapat digunakan di sini.
\end{example}

Uji akar atau uji rasio sebagaimana diberikan tidak selalu dapat diterapkan. Artinya,
limit dari
$\bigl \lvert \frac{a_{k+1}}{a_k} \bigr \rvert$
atau
$\sqrt[k]{\lvert a_k \rvert}$
mungkin tidak ada.
Kadang-kadang uji tersebut harus diterapkan pada deret itu sendiri---untuk mencari
$L$ alih-alih $A$. Sebagai contoh, untuk deret
$\sum_{k=1}^\infty 2^k x^{2k}$, kita tidak dapat menerapkan uji akar dengan $A$; kita
harus bekerja dengan limit $\sqrt[k]{\lvert 2^k x^{2k} \rvert}$,
yang sama dengan $L = 2 {\lvert x \rvert}^2$. Nilai $L$ ini kurang dari $1$ ketika
$\lvert x \rvert < \frac{1}{\sqrt{2}}$, sehingga $\frac{1}{\sqrt{2}}$
adalah jari-jari konvergensinya.
Terdapat cara lain untuk mencari jari-jari tersebut, tetapi uji rasio dan uji akar
mencakup banyak deret yang muncul dalam praktik.
%Terdapat cara-cara yang lebih canggih untuk mencari jari-jari konvergensi,
%tetapi cara-cara tersebut akan berada di luar cakupan bab ini.
%Kedua metode di atas
%mencakup banyak deret yang muncul dalam praktik.
%Sering kali kedua uji dapat diterapkan, meskipun limitnya mungkin
%lebih mudah dihitung dengan salah satu uji daripada yang lain.

Di titik-titik ujung,
$x = x_0-\rho$ dan $x = x_0+\rho$,
deret tersebut mungkin konvergen atau mungkin tidak, dan uji-uji di atas
tidak menyatakan apa pun tentang konvergensi di sana.
Kadang-kadang konvergensi di titik ujung penting, tetapi untuk
keperluan kita, kita tidak akan terlalu mengkhawatirkannya.

\subsection{Fungsi analitik}

Fungsi-fungsi yang direpresentasikan oleh deret pangkat disebut
\emph{\myindex{fungsi analitik}}. Tidak semua fungsi analitik,
tetapi kemungkinan besar semua fungsi yang telah Anda lihat dalam kalkulus bersifat analitik.
Suatu fungsi analitik $f(x)$ sama dengan \emph{\myindex{deret Taylor}}-nya%
\footnote{Dinamai menurut matematikawan Inggris
\href{http://en.wikipedia.org/wiki/Brook_Taylor}{Sir Brook Taylor}
(1685--1731).} (suatu deret pangkat yang dihitung dari $f$)
untuk $x$ di dekat suatu titik $x_0$ yang diberikan:
\begin{equation} \label{ps:tayloreq}
f(x) = \sum_{k=0}^\infty \frac{f^{(k)}(x_0)}{k!} {(x-x_0)}^k 
=
f(x_0)
+ f'(x_0) (x-x_0)
+ \frac{f''(x_0)}{2} (x-x_0)^2
+ \cdots
,
\avoidbreak
\end{equation}
di mana $f^{(k)}(x_0)$ menyatakan turunan berorde $k^{\text{th}}$ dari $f(x)$
di titik $x_0$.

Sebagai contoh, sinus adalah fungsi analitik dan deret Taylornya
di sekitar $x_0 = 0$
diberikan oleh
\begin{equation*}
\sin(x) = \sum_{n=0}^\infty \frac{{(-1)}^n}{(2n+1)!}
 x^{2n+1} .
\end{equation*}
Dalam \figurevref{ps:sin}, kita menggambar $\sin(x)$ dan pemotongan
deret hingga derajat 5 dan 9. Anda dapat melihat bahwa hampiran tersebut sangat
baik untuk $x$ di dekat 0, tetapi memburuk untuk $x$ yang lebih jauh dari 0. Inilah 
yang terjadi secara umum.
Untuk memperoleh hampiran yang baik jauh dari $x_0$, Anda
perlu mengambil semakin banyak suku dari deret Taylor.

\begin{myfig}
\capstart
\diffyincludegraphics{width=3in}{width=4.5in}{ps-sin}{%
Grafik fungsi sinus pada interval minus 10 hingga 10.
Kemudian digambar dua hampiran yang keduanya muncul dari bagian bawah
gambar dan mendekati grafik fungsi sinus; keduanya sangat dekat dengan grafik tersebut
dalam suatu interval di sekitar titik asal, lalu setelah beberapa saat keduanya meninggalkan grafik melalui
bagian atas gambar. Salah satu hampiran dekat dengan
grafik sinus pada interval yang lebih panjang.}
\caption{Fungsi sinus dan hampiran Taylornya
di sekitar $x_0=0$
berderajat $5^{\text{th}}$ dan $9^{\text{th}}$.\label{ps:sin}}
\end{myfig}

\subsection{Memanipulasi deret pangkat}

Salah satu sifat utama deret pangkat yang akan kita gunakan adalah
bahwa kita dapat mendiferensialkannya suku demi suku. Artinya,
andaikan bahwa 
$\sum a_k {(x-x_0)}^k$ adalah deret pangkat konvergen. Maka
untuk $x$ di dalam jari-jari konvergensi, kita mempunyai
\begin{equation*}
\frac{d}{dx}
\left[\sum_{k=0}^\infty a_k {(x-x_0)}^k\right]
=
\sum_{k=1}^\infty k a_k {(x-x_0)}^{k-1}
=
a_1
+ 2 a_2 (x-x_0)
+ 3 a_3 {(x-x_0)}^{2}
+ \cdots
.
\end{equation*}
Perhatikan bahwa suku yang bersesuaian dengan $k=0$ lenyap karena
suku tersebut konstan. Jari-jari konvergensi deret yang telah didiferensialkan
sama dengan jari-jari deret asal.

\begin{example}
Mari kita tunjukkan bahwa fungsi eksponensial $y=e^x$ memenuhi $y'=y$. Andaikan
kita belum mengetahui hal itu.
Tuliskan
\begin{equation*}
y = e^x = \sum_{k=0}^\infty \frac{1}{k!} x^k .
\end{equation*}
Diferensialkan $y$,
\begin{equation*}
y' = \sum_{k=1}^\infty k \frac{1}{k!} x^{k-1} =
\sum_{k=1}^\infty \frac{1}{(k-1)!} x^{k-1} .
\end{equation*}
Kita \emph{mengindeks ulang}\index{pengindeksan ulang deret}
deret tersebut hanya dengan mengganti $k$ dengan $k+1$. Deret tersebut
tidak berubah; yang berubah hanyalah cara kita menuliskannya. Setelah
pengindeksan ulang, deret kembali dimulai
pada $k=0$.
\begin{equation*}
\sum_{k=1}^\infty \frac{1}{(k-1)!} x^{k-1} =
\sum_{k+1=1}^\infty \frac{1}{\bigl((k+1)-1\bigr)!} x^{(k+1)-1} =
\sum_{k=0}^\infty \frac{1}{k!} x^k .
\end{equation*}
Itu tepat merupakan deret pangkat untuk $e^x$ yang kita gunakan di awal,
sehingga kita telah menunjukkan bahwa $\frac{d}{dx} [ e^x ] = e^x$.
\end{example}

Deret pangkat konvergen dapat dijumlahkan dan dikalikan satu sama lain, serta dikalikan
dengan konstanta menggunakan aturan-aturan berikut. Pertama, kita dapat menjumlahkan deret dengan
menjumlahkan suku demi suku,
\begin{equation*}
\left(\sum_{k=0}^\infty a_k {(x-x_0)}^k\right)
+
\left(\sum_{k=0}^\infty b_k {(x-x_0)}^k\right)
=
\sum_{k=0}^\infty (a_k+b_k) {(x-x_0)}^k .
\end{equation*}
Kita dapat mengalikan dengan konstanta,
\begin{equation*}
\alpha
\left(\sum_{k=0}^\infty a_k {(x-x_0)}^k\right)
=
\sum_{k=0}^\infty \alpha a_k {(x-x_0)}^k .
\end{equation*}
Kita juga dapat mengalikan deret satu sama lain,
\begin{equation*}
\left(\sum_{k=0}^\infty a_k {(x-x_0)}^k\right)
\,
\left(\sum_{k=0}^\infty b_k {(x-x_0)}^k\right)
=
\sum_{k=0}^\infty c_k {(x-x_0)}^k ,
\end{equation*}
di mana
$c_k = a_0b_k + a_1 b_{k-1} + \cdots + a_k b_0$.
Jari-jari konvergensi jumlah atau hasil kali
setidaknya sama dengan minimum dari jari-jari konvergensi
kedua deret yang terlibat.

\subsection{Deret pangkat untuk fungsi rasional}

Polinomial hanyalah deret pangkat berhingga. Artinya, suatu polinomial
adalah deret pangkat dengan
$a_k$ bernilai nol untuk semua $k$ yang cukup besar. Kita selalu dapat mengembangkan
suatu polinomial sebagai deret pangkat di sekitar sembarang titik $x_0$ dengan menuliskan
polinomial tersebut sebagai polinomial dalam $(x-x_0)$. Sebagai contoh,
mari kita tuliskan
$2x^2-3x+4$ sebagai deret pangkat di sekitar $x_0 = 1$:
\begin{equation*}
2x^2-3x+4 = 3 + (x-1) + 2{(x-1)}^2 .
\end{equation*}
Dengan kata lain, $a_0 = 3$, $a_1 = 1$, $a_2 = 2$, dan semua
$a_k = 0$ lainnya. Untuk melakukan ini, kita mengetahui bahwa $a_k = 0$ untuk semua $k \geq 3$ karena
polinomial tersebut berderajat 2.
Kita menulis $a_0 + a_1(x-1) + a_2{(x-1)}^2$, mengembangkannya, lalu menyelesaikan
untuk $a_0$, $a_1$, dan $a_2$. Kita juga dapat mendiferensialkan di $x=1$
dan menggunakan rumus deret Taylor \eqref{ps:tayloreq}.

Mari kita tinjau fungsi rasional, yaitu hasil bagi polinomial.
Suatu fakta penting adalah 
bahwa suatu deret untuk fungsi hanya mendefinisikan fungsi tersebut
pada suatu interval, meskipun fungsi itu didefinisikan di tempat lain. Sebagai contoh, untuk
$-1 < x < 1$,
\begin{equation*}
\frac{1}{1-x} =
\sum_{k=0}^\infty x^k =
1 + x + x^2 + \cdots
\end{equation*}
Deret ini disebut \emph{\myindex{deret geometri}}. Uji rasio
memberi tahu kita bahwa jari-jari konvergensinya adalah $1$. Deret tersebut
divergen untuk $x \leq -1$ dan $x \geq 1$, meskipun
$\frac{1}{1-x}$ didefinisikan untuk semua $x \neq 1$.

Alih-alih menerapkan
pengembangan deret Taylor \eqref{ps:tayloreq},
deret geometri beserta aturan penjumlahan dan
perkalian deret pangkat dapat digunakan untuk mengembangkan sembarang fungsi rasional di sekitar
suatu titik $x_0$, selama penyebutnya tidak nol di $x_0$.

\begin{example}
Kembangkan $\frac{x}{1+2x+x^2}$ sebagai deret pangkat di sekitar titik asal ($x_0 = 0$) dan
carilah jari-jari konvergensinya.

Pertama, tuliskan $1+2x+x^2 = {(1+x)}^2 = {\bigl(1-(-x)\bigr)}^2$.
Hitung
\begin{equation*}
\begin{split}
\frac{x}{1+2x+x^2}
&=
x \,
{\left(
\frac{1}{1-(-x)}
\right)}^2
\\
&=
x \,
{ \left( 
\sum_{k=0}^\infty {(-1)}^k x^k 
\right)}^2
\\
&=
x \,
\left(
\sum_{k=0}^\infty c_k x^k 
\right)
\\
&=
\sum_{k=0}^\infty c_k x^{k+1} ,
\end{split}
\end{equation*}
di mana untuk memperoleh $c_k$, kita menggunakan rumus hasil kali deret:
$c_0 = 1$, $c_1 = -1 -1 = -2$, $c_2 = 1+1+1 = 3$, dan seterusnya.
Oleh karena itu,
\begin{equation*}
\frac{x}{1+2x+x^2}
=
\sum_{k=1}^\infty {(-1)}^{k+1} k x^k
= x-2x^2+3x^3-4x^4+\cdots
\end{equation*}
Jari-jari konvergensinya setidaknya 1. Kita menggunakan uji rasio
\begin{equation*}
\lim_{k\to\infty}
\left\lvert \frac{a_{k+1}}{a_k} \right\rvert
=
\lim_{k\to\infty}
\left\lvert \frac{{(-1)}^{k+2} (k+1)}{{(-1)}^{k+1}k} \right\rvert
=
\lim_{k\to\infty}
\frac{k+1}{k}
= 1 .
\end{equation*}
Jadi jari-jari konvergensinya benar-benar sama dengan 1.
\end{example}

Jika fungsi rasionalnya lebih rumit, kita juga dapat
menggunakan metode pecahan parsial. Sebagai contoh,
untuk mencari deret Taylor bagi $\frac{x^3+x}{x^2-1}$, kita menulis
\begin{equation*}
\frac{x^3+x}{x^2-1}
=
x + \frac{1}{1+x} - \frac{1}{1-x}
=
x + \sum_{k=0}^\infty {(-1)}^k x^k - \sum_{k=0}^\infty x^k
=
- x + \sum_{\substack{k=3 \\ k \text{ ganjil}}}^\infty (-2) x^k .
\end{equation*}

\subsection{Latihan}

\begin{exercise}
Apakah $\displaystyle \sum_{k=0}^\infty e^k x^k$ merupakan deret pangkat konvergen?
Jika ya, carilah jari-jari konvergensinya.
\end{exercise}

\begin{exercise}
Apakah $\displaystyle \sum_{k=0}^\infty k x^k$ merupakan deret pangkat konvergen?
Jika ya, carilah jari-jari konvergensinya.
\end{exercise}

\begin{exercise}
Apakah $\displaystyle \sum_{k=0}^\infty k! x^k$ merupakan deret pangkat konvergen?
Jika ya, carilah jari-jari konvergensinya.
\end{exercise}

\begin{exercise}
Apakah $\displaystyle \sum_{k=0}^\infty \frac{1}{(2k)!} {(x-10)}^k$
merupakan deret pangkat konvergen? Jika ya, carilah jari-jari konvergensinya.
\end{exercise}

\begin{exercise}
Tentukan deret Taylor untuk $\sin x$ di sekitar titik $x_0 = \pi$.
\end{exercise}

\begin{exercise}
Carilah deret Taylor untuk $\ln x$ di sekitar titik $x_0 = 1$,
beserta jari-jari konvergensinya.
\end{exercise}

\begin{exercise}
Tentukan deret Taylor
beserta jari-jari konvergensinya untuk $\dfrac{1}{1+x}$
di sekitar $x_0 = 0$.
\end{exercise}

\begin{exercise}
Tentukan deret Taylor beserta jari-jari konvergensinya
untuk
$\dfrac{x}{4-x^2}$ di sekitar $x_0 = 0$. Petunjuk: Anda tidak akan dapat
menggunakan uji rasio.
\end{exercise}

\begin{exercise}
Kembangkan $x^5+5x+1$ sebagai deret pangkat di sekitar $x_0 = 5$.
\end{exercise}

\begin{exercise}
Andaikan bahwa uji rasio dapat diterapkan pada suatu deret
$\displaystyle \sum_{k=0}^\infty a_k x^k$. Tunjukkan, dengan menggunakan uji
rasio, bahwa jari-jari konvergensi deret yang didiferensialkan
sama dengan jari-jari konvergensi deret asal.
\end{exercise}

\begin{exercise}
Andaikan $f$ adalah fungsi analitik sedemikian sehingga
$f^{(n)}(0) = n$. Carilah $f(1)$.
\end{exercise}

\setcounter{exercise}{100}

\begin{exercise}
Apakah
$\displaystyle \sum_{n=1}^\infty {(0.1)}^n x^n$
merupakan deret pangkat konvergen?
Jika ya, carilah jari-jari konvergensinya.
\end{exercise}
\exsol{%
Ya. Jari-jari konvergensinya adalah $10$.
}

\begin{exercise}[menantang]
Apakah
$\displaystyle \sum_{n=1}^\infty \frac{n!}{n^n} x^n$
merupakan deret pangkat konvergen?
Jika ya, carilah jari-jari konvergensinya.
\end{exercise}
\exsol{%
Ya. Jari-jari konvergensinya adalah $e$.
}

\begin{exercise}
Dengan menggunakan deret geometri, kembangkan $\frac{1}{1-x}$ di sekitar $x_0=2$.
Untuk nilai $x$ apa deret tersebut konvergen?
\end{exercise}
\exsol{%
$\frac{1}{1-x} = -\frac{1}{1-(2-x)}$, sehingga
$\frac{1}{1-x} =
\sum\limits_{n=0}^\infty {(-1)}^{n+1} {(x-2)}^n$,
yang konvergen untuk $1 < x < 3$.
}

\begin{exercise}[menantang]
Carilah deret Taylor untuk $x^7 e^x$ di sekitar $x_0 = 0$.
\end{exercise}
\exsol{%
$\sum\limits_{n=7}^\infty
\frac{1}{(n-7)!} x^n$
}

\begin{exercise}[menantang]
Bayangkan $f$ dan $g$ adalah fungsi analitik sedemikian sehingga
$f^{(k)}(0) = g^{(k)}(0)$ untuk semua $k$ yang cukup besar. Apa yang dapat Anda
katakan tentang $f(x)-g(x)$?
\end{exercise}
\exsol{%
$f(x)-g(x)$ adalah suatu polinomial. Petunjuk: Gunakan deret Taylor.
}

\begin{exercise}
Indeks ulang deret-deret berikut agar pangkatnya berbentuk $x^k$.
\begin{tasks}(3)
\task $\displaystyle \sum_{k=3}^\infty k(k-1)x^{k+3}$
\task $\displaystyle \sum_{k=1}^\infty (k+1)x^{k-1}$
\task $\displaystyle \sum_{k=5}^\infty 2kx^{k-2}$
\end{tasks}
\end{exercise}
\exsol{%
a) $\displaystyle \sum_{k=6}^\infty (k-3)(k-4)x^{k}$
\quad
b) $\displaystyle \sum_{k=0}^\infty (k+2)x^{k}$
\quad
c) $\displaystyle \sum_{k=3}^\infty 2(k+2)x^{k}$
}


%%%%%%%%%%%%%%%%%%%%%%%%%%%%%%%%%%%%%%%%%%%%%%%%%%%%%%%%%%%%%%%%%%%%%%%%%%%%%%

\sectionnewpage
\section{Solusi deret untuk PDB linear orde dua}
\label{seriessols:section}

\sectionnotes{1,5 atau 2 kuliah\EPref{, \S8.2 dalam \cite{EP}}\BDref{,
\S5.2 dan \S5.3 dalam \cite{BD}}}

Tinjau suatu PDB linear homogen orde dua berbentuk
\begin{equation*}
P(x) y'' + Q(x) y' + R(x) y = 0 .
\end{equation*}
Andaikan $P(x)$, $Q(x)$, dan $R(x)$ adalah polinomial. Kita akan 
mencoba solusi berbentuk
\begin{equation*}
y = \sum_{k=0}^\infty a_k {(x-x_0)}^k ,
\end{equation*}
dan menyelesaikan untuk $a_k$ guna memperoleh solusi yang didefinisikan pada suatu
interval di sekitar $x_0$.

Titik $x_0$ disebut \emph{\myindex{titik biasa}}
jika $P(x_0) \neq 0$. Artinya, fungsi-fungsi
\begin{equation*}
\frac{Q(x)}{P(x)} \qquad \text{dan} \qquad \frac{R(x)}{P(x)}
\end{equation*}
didefinisikan untuk $x$ di dekat $x_0$. Jika $P(x_0) = 0$, maka $x_0$
adalah suatu \emph{\myindex{titik singular}}. Menangani titik singular
lebih sulit daripada menangani titik biasa, sehingga sekarang kita hanya berfokus pada titik biasa.

\begin{example}
Kita mulai dengan contoh yang sangat sederhana
\begin{equation*}
y'' - y = 0 .
\end{equation*}
Mari kita coba solusi deret pangkat di dekat $x_0 = 0$,
yang merupakan titik biasa. Bahkan, setiap titik adalah titik
biasa karena persamaan tersebut berkoefisien konstan. Kita sudah tahu
bahwa kita seharusnya memperoleh fungsi eksponensial atau sinus dan kosinus hiperbolik,
tetapi mari kita berpura-pura belum mengetahui fakta ini.

Kita coba
\begin{equation*}
y = \sum_{k=0}^\infty a_k x^k .
\end{equation*}
Jika kita mendiferensialkan, suku $k=0$ adalah konstanta sehingga lenyap.
Kita memperoleh
\begin{equation*}
y' = \sum_{k=1}^\infty k a_k x^{k-1} .
\end{equation*}
Kita mendiferensialkan sekali lagi untuk memperoleh (sekarang suku $k=1$ lenyap)
\begin{equation*}
y'' = \sum_{k=2}^\infty k(k-1) a_k x^{k-2} .
\end{equation*}
Kita perlu mengurangkan $y$ dari $y''$, sehingga 
kita mengindeks ulang deret untuk $y''$ (mengganti $k$ dengan $k+2$):
\begin{equation*}
y'' = \sum_{k=0}^\infty (k+2)\,(k+1) \, a_{k+2} x^k .
\end{equation*}
Sekarang kita substitusikan $y$ dan $y''$ ke dalam persamaan diferensial
\begin{equation*}
\begin{split}
0 = y''-y & = 
\Biggl( \sum_{k=0}^\infty (k+2)\,(k+1) \, a_{k+2} x^k  \Biggr)
-
\Biggl( \sum_{k=0}^\infty a_k x^k \Biggr)
\\
& =
\sum_{k=0}^\infty \,\Bigl( (k+2)\,(k+1) \, a_{k+2} x^k 
-
a_k x^k \Bigr)
\\
& =
\sum_{k=0}^\infty \,\bigl( (k+2)\,(k+1) \,a_{k+2} - a_k \bigr) \, x^k  .
\end{split}
\end{equation*}
Karena $y'' - y$ seharusnya sama dengan 0, kita tahu bahwa
koefisien-koefisien deret yang dihasilkan harus sama dengan 0. Oleh karena itu,
untuk setiap $k=0,1,2,3,\ldots$, kita mempunyai
\begin{equation*}
(k+2)\,(k+1) \,a_{k+2} - a_k = 0 ,
\qquad
\text{atau}
\qquad
a_{k+2} = \frac{a_k}{(k+2)(k+1)} .
\end{equation*}
Persamaan di atas disebut \emph{\myindex{relasi rekurensi}}
untuk koefisien-koefisien deret pangkat.
Nilai $a_0$ atau $a_1$ tidak menjadi masalah. Keduanya dapat dipilih sembarang.
Namun, setelah kita memilih $a_0$ dan $a_1$, semua koefisien lainnya
ditentukan oleh relasi rekurensi.

Mari kita lihat nilai yang harus dimiliki
koefisien-koefisien tersebut. Pertama, $a_0$ dan $a_1$ sembarang. Kemudian,
\begin{equation*}
a_2 = \frac{a_0}{2}, \quad
a_3 = \frac{a_1}{3 \cdot 2}, \quad
a_4 = \frac{a_2}{4 \cdot 3} = \frac{a_0}{4 \cdot 3 \cdot 2}, \quad
a_5 = \frac{a_3}{5 \cdot 4} = \frac{a_1}{5 \cdot 4 \cdot 3 \cdot 2}, \quad \ldots
\end{equation*}
Untuk $k$ genap, yaitu $k=2n$,
kita mempunyai
\begin{equation*}
a_k = a_{2n} = \frac{a_0}{(2n)!} .
\end{equation*}
Untuk $k$ ganjil, yaitu $k=2n+1$, kita mempunyai
\begin{equation*}
a_k = a_{2n+1} = \frac{a_1}{(2n+1)!} .
\end{equation*}
Kita tuliskan deret untuk solusinya
\begin{equation*}
%mbxlatex \begin{split}
y
%mbxlatex &
=
\sum_{k=0}^\infty
a_k x^k
=
\sum_{n=0}^\infty
\left(
\frac{a_0}{(2n)!} \,x^{2n}
+
\frac{a_1}{(2n+1)!} \,x^{2n+1}
\right)
%mbxlatex \\
%mbxlatex &
=
a_0
\sum_{n=0}^\infty
\frac{1}{(2n)!} \,x^{2n}
+
a_1
\sum_{n=0}^\infty
\frac{1}{(2n+1)!} \,x^{2n+1} .
%mbxlatex \end{split}
\end{equation*}
Kita mengenali kedua deret tersebut sebagai kosinus dan sinus hiperbolik.
Oleh karena itu,
\begin{equation*}
y =
a_0 \cosh x + a_1 \sinh x .
\end{equation*}
\end{example}

Tentu saja, secara umum kita tidak akan dapat mengenali 
deret yang muncul, karena biasanya tidak akan ada
fungsi elementer yang cocok dengannya. Dalam hal itu, kita akan
puas dengan bentuk deretnya.

\begin{example}
Mari kita kerjakan contoh yang lebih rumit. Tinjau
\emph{\myindex{persamaan Airy}}%
\footnote{Dinamai menurut matematikawan Inggris
\href{http://en.wikipedia.org/wiki/George_Biddell_Airy}{Sir George Biddell Airy}
(1801--1892).}:
\begin{equation*}
y'' - xy = 0 ,
\end{equation*}
di dekat titik $x_0 = 0$. Perhatikan bahwa $x_0 = 0$ adalah titik biasa.

\pagebreak[2]
Kita coba
\begin{equation*}
y = \sum_{k=0}^\infty a_k x^k .
\end{equation*}
Kita mendiferensialkan dua kali (seperti di atas) untuk memperoleh
\begin{equation*}
y'' = \sum_{k=2}^\infty k\,(k-1) \, a_k x^{k-2} .
\end{equation*}
Kita substitusikan $y$ ke dalam persamaan:
\begin{equation*}
\begin{split}
0 = y''-xy &= 
\Biggl( \sum_{k=2}^\infty k\,(k-1) \, a_k x^{k-2}  \Biggr)
-
x
\Biggl( \sum_{k=0}^\infty a_k x^k \Biggr)
\\
&=
\Biggl( \sum_{k=2}^\infty k\,(k-1) \, a_k x^{k-2}  \Biggr)
-
\Biggl( \sum_{k=0}^\infty a_k x^{k+1} \Biggr) .
\end{split}
\end{equation*}
Kita mengindeks ulang dan menyelaraskan indeks awal agar deret-deret tersebut
dapat dijumlahkan:
\begin{equation*}
\begin{split}
0 = y''-xy
&= 
\Biggl( 2 a_2 + \sum_{k=1}^\infty (k+2)\,(k+1) \, a_{k+2} x^k  \Biggr)
-
\Biggl( \sum_{k=1}^\infty a_{k-1} x^k \Biggr)
\\
&= 
2 a_2 + 
\sum_{k=1}^\infty \Bigl( (k+2)\,(k+1) \, a_{k+2} - a_{k-1} \Bigr) \, x^k .
\end{split}
\end{equation*}
Karena $y''-xy$ seharusnya bernilai $0$, kita mempunyai $a_2 = 0$, dan
\begin{equation*}
(k+2)\,(k+1) \,a_{k+2} - a_{k-1} = 0 ,
\qquad
\text{atau}
\qquad
a_{k+2} = \frac{a_{k-1}}{(k+2)(k+1)} .
\end{equation*}
Kita mempunyai konstanta sembarang $a_0$ dan $a_1$, serta $a_2 = 0$.
Kita menghitung beberapa koefisien yang lebih tinggi:
\begin{multline*}
a_3 = \frac{a_0}{3 \cdot 2}, \quad
a_4 = \frac{a_1}{4 \cdot 3}, \quad
a_5 = \frac{a_2}{5 \cdot 4} = 0, \quad
a_6 = \frac{a_3}{6 \cdot 5} = \frac{a_0}{6 \cdot 5 \cdot 3 \cdot 2},
\\
a_7 = \frac{a_4}{7 \cdot 6} = \frac{a_1}{7 \cdot 6 \cdot 4 \cdot 3}, \quad
a_8 = \frac{a_5}{8 \cdot 7} = 0, \quad
a_9 = \frac{a_6}{9 \cdot 8} = \frac{a_0}{9 \cdot 8 \cdot 6 \cdot 5 \cdot 3
\cdot 2}, \quad
\ldots
\end{multline*}
Kita melangkah dalam kelipatan tiga.
Untuk $a_k$ dengan $k$ kelipatan $3$, yaitu $k=3n$, kita melihat
bahwa
\begin{equation*}
a_{3n} = \frac{a_0}{2 \cdot 3 \cdot 5 \cdot 6 \cdots (3n-1) \cdot (3n)} .
\end{equation*}
Untuk $a_k$ dengan $k = 3n+1$, kita melihat
\begin{equation*}
a_{3n+1} = \frac{a_1}{3 \cdot 4 \cdot 6 \cdot 7 \cdots (3n) \cdot (3n+1)} .
\end{equation*}
Terakhir, karena $a_2 = 0$, kita mempunyai $a_5 = 0$, $a_8 = 0$, $a_{11}=0$, dan seterusnya.
Secara umum, $a_{3n+2} = 0$.

Dengan kata lain, jika kita menuliskan deret untuk $y$,
deret tersebut mempunyai dua bagian
\begin{equation*}
\begin{split}
y &=
\left(
a_0 + \frac{a_0}{6} x^3 + \frac{a_0}{180} x^6 + \cdots +
\frac{a_0}{2 \cdot 3 \cdot 5 \cdot 6 \cdots (3n-1) \cdot (3n)} x^{3n} + \cdots
\right)
\\
&\phantom{=}
+
\left(
a_1 x + \frac{a_1}{12} x^4 + \frac{a_1}{504} x^7 + \cdots +
\frac{a_1}{3 \cdot 4 \cdot 6 \cdot 7 \cdots (3n) \cdot (3n+1)} x^{3n+1} + \cdots
\right)
\\
& =
a_0
\left(
1 + \frac{1}{6} x^3 + \frac{1}{180} x^6 + \cdots +
\frac{1}{2 \cdot 3 \cdot 5 \cdot 6 \cdots (3n-1) \cdot (3n)} x^{3n} + \cdots
\right)
\\
&\phantom{=}
+
a_1
\left(
x + \frac{1}{12} x^4 + \frac{1}{504} x^7 + \cdots +
\frac{1}{3 \cdot 4 \cdot 6 \cdot 7 \cdots (3n) \cdot (3n+1)} x^{3n+1} + \cdots
\right) .
\end{split}
\end{equation*}
Kita definisikan
\begin{align*}
y_1(x) &= 
1 + \frac{1}{6} x^3 + \frac{1}{180} x^6 + \cdots +
\frac{1}{2 \cdot 3 \cdot 5 \cdot 6 \cdots (3n-1) \cdot (3n)} x^{3n} + \cdots, \\
y_2(x) &= 
x + \frac{1}{12} x^4 + \frac{1}{504} x^7 + \cdots +
\frac{1}{3 \cdot 4 \cdot 6 \cdot 7 \cdots (3n) \cdot (3n+1)} x^{3n+1} + \cdots ,
\end{align*}
dan tuliskan solusi umum persamaan tersebut sebagai
$y(x)= a_0 y_1(x) + a_1 y_2(x)$. Jika kita mensubstitusikan $x=0$ ke dalam
deret pangkat untuk $y_1$ dan $y_2$, kita memperoleh
$y_1(0) = 1$ dan $y_2(0) = 0$. Demikian pula,
$y_1'(0) = 0$ dan $y_2'(0) = 1$. Oleh karena itu, $y = a_0 y_1 + a_1 y_2$
adalah suatu solusi
yang memenuhi syarat awal $y(0) = a_0$ dan $y'(0) = a_1$.

\begin{myfig}
\capstart
\diffyincludegraphics{width=3in}{width=4.5in}{ps-airy}{%
Grafik dua fungsi; keduanya tampak berosilasi naik-turun
untuk x kurang dari 0, dan keduanya tampak naik seperti fungsi eksponensial untuk
x positif. Salah satunya melalui titik asal dengan kemiringan 1, sedangkan yang lain
melalui titik (0,1) dengan kemiringan nol.}
\caption{Dua solusi $y_1$ dan $y_2$ untuk persamaan Airy.\label{ps:airyfig}}
\end{myfig}
\end{example}
Fungsi-fungsi $y_1$ dan $y_2$ tidak dapat dituliskan dalam bentuk fungsi
elementer yang Anda ketahui. Lihat \figurevref{ps:airyfig} untuk grafik
solusi $y_1$ dan $y_2$. Fungsi-fungsi ini mempunyai banyak sifat menarik.
Sebagai contoh, keduanya berosilasi untuk $x$ negatif
(seperti solusi $y''+y=0$) dan
untuk $x$ positif keduanya tumbuh tanpa batas (seperti solusi $y''-y=0$).

\medskip

Kadang-kadang suatu solusi deret ternyata merupakan polinomial.
Mari kita lihat suatu contoh.

\begin{example}
Mari kita cari solusi untuk apa yang disebut
\emph{\myindex{persamaan Hermite berorde $n$}}%
\footnote{Dinamai menurut matematikawan Prancis
\href{http://en.wikipedia.org/wiki/Hermite}{Charles Hermite}
(1822--1901).}:
\begin{equation*}
y'' -2xy' + 2n y = 0 .
\end{equation*}

Mari kita cari solusi di sekitar titik $x_0 = 0$.
Kita coba
\begin{equation*}
y = \sum_{k=0}^\infty a_k x^k .
\end{equation*}
Kita mendiferensialkan (seperti di atas) untuk memperoleh
\begin{equation*}
y' = \sum_{k=1}^\infty k a_k x^{k-1} 
\qquad \text{dan} \qquad
y'' = \sum_{k=2}^\infty k\,(k-1) \, a_k x^{k-2} .
\end{equation*}

Sekarang kita substitusikan ke dalam persamaan
\begin{equation*}
\begin{split}
0 &= y''-2xy'+2ny \\
 &= 
\Biggl( \sum_{k=2}^\infty k(k-1) a_k x^{k-2}  \Biggr)
-
2x
\Biggl( \sum_{k=1}^\infty k a_k x^{k-1} \Biggr)
+
2n
\Biggl( \sum_{k=0}^\infty a_k x^k \Biggr)
\\
&=
\Biggl( \sum_{k=2}^\infty k(k-1) a_k x^{k-2}  \Biggr)
-
\Biggl( \sum_{k=1}^\infty 2k a_k x^k \Biggr)
+
\Biggl( \sum_{k=0}^\infty 2n a_k x^k \Biggr)
\\
&=
\Biggl(2a_2+
 \sum_{k=1}^\infty (k+2)(k+1) a_{k+2} x^k  \Biggr)
-
\Biggl( \sum_{k=1}^\infty 2k a_k x^k \Biggr)
+
\Biggl(
2na_0 + 
\sum_{k=1}^\infty 2n a_k x^k \Biggr)
\\
&=
2a_2+2na_0+
\sum_{k=1}^\infty \bigl( (k+2)(k+1)  a_{k+2} - 2ka_k + 2n a_k \bigr) x^k .
\end{split}
\end{equation*}
Karena $y''-2xy'+2ny = 0$, kita mempunyai
\begin{equation*}
(k+2)(k+1)  a_{k+2} + ( - 2k+ 2n) a_k = 0 ,
\qquad
\text{atau}
\qquad
a_{k+2} = \frac{(2k-2n)}{(k+2)(k+1)} a_k .
\end{equation*}
Relasi rekurensi ini sebenarnya juga mencakup
$a_2 = -na_0$ (yang berasal dari suku konstan $2a_2+2na_0 = 0$).
Sekali lagi, $a_0$ dan $a_1$ sembarang.
\begin{align*}
& a_2 = \frac{-2n}{2 \cdot 1}a_0,
\displaybreak[0]\\
& a_3 = \frac{2(1-n)}{3 \cdot 2} a_1,
\displaybreak[0]\\
& a_4 = \frac{2(2-n)}{4 \cdot 3} a_2
=
\frac{2^2(2-n)(-n)}{4 \cdot 3 \cdot 2 \cdot 1} a_0 ,
\displaybreak[0]\\
&
a_5 = \frac{2(3-n)}{5 \cdot 4} a_3
=
\frac{2^2(3-n)(1-n)}{5 \cdot 4 \cdot 3 \cdot 2} a_1 ,
\quad \ldots
\end{align*}
Kita pisahkan koefisien genap dan ganjil untuk memperoleh
\begin{align*}
a_{2m} &=\frac{2^m(-n)(2-n)\cdots(2m-2-n)}{(2m)!} a_0, \\
a_{2m+1} &=\frac{2^m(1-n)(3-n)\cdots(2m-1-n)}{(2m+1)!} a_1 .
\end{align*}
Kita tuliskan kedua deret tersebut, satu dengan pangkat genap dan satu dengan pangkat
ganjil:
\begin{align*}
y_1(x) & = 
1+\frac{2(-n)}{2!} x^2 + \frac{2^2(-n)(2-n)}{4!} x^4 + 
\frac{2^3(-n)(2-n)(4-n)}{6!} x^6 + \cdots ,
\\
y_2(x) & = 
x+\frac{2(1-n)}{3!} x^3 + \frac{2^2(1-n)(3-n)}{5!} x^5 + 
\frac{2^3(1-n)(3-n)(5-n)}{7!} x^7 + \cdots .
\end{align*}
Kemudian
\begin{equation*}
y(x) = a_0 y_1(x) + a_1 y_2(x) .
\end{equation*}

Kita mencatat bahwa jika $n$ adalah bilangan bulat positif genap, maka $y_1(x)$ adalah
polinomial karena semua koefisien dalam deret setelah
derajat $n$ bernilai nol. Jika $n$ adalah bilangan bulat positif ganjil, maka $y_2(x)$ adalah
polinomial. Sebagai contoh, jika $n=4$, maka
\begin{equation*}
y_1(x) = 1 + \frac{2(-4)}{2!} x^2 + \frac{2^2(-4)(2-4)}{4!} x^4
= 1 - 4x^2 + \frac{4}{3} x^4 .
\end{equation*}
\end{example}

\subsection{Latihan}

Dalam latihan-latihan berikut, ketika diminta menyelesaikan suatu persamaan menggunakan metode
deret pangkat, Anda harus mencari beberapa suku pertama deret tersebut,
dan jika mungkin mencari rumus umum untuk koefisien berindeks $k^{\text{th}}$.

\begin{exercise}
Gunakan metode deret pangkat untuk menyelesaikan $y''+y = 0$ di titik $x_0 = 1$.
\end{exercise}

\begin{exercise}
Gunakan metode deret pangkat untuk menyelesaikan $y''+4xy = 0$ di titik $x_0 = 0$.
\end{exercise}

\begin{exercise}
Gunakan metode deret pangkat untuk menyelesaikan $y''-xy = 0$ di titik $x_0 = 1$.
\end{exercise}

\begin{exercise}
Gunakan metode deret pangkat untuk menyelesaikan $y''+x^2y = 0$ di titik $x_0 = 0$.
\end{exercise}

\begin{exercise}
Metode-metode tersebut juga berlaku untuk orde selain orde dua. Cobalah metode
dalam bagian ini untuk menyelesaikan persamaan orde pertama $y'-xy = 0$ di
titik $x_0 = 0$.
\end{exercise}

\begin{exercise}[\myindex{persamaan Chebyshev berorde $p$}]
\leavevmode
\begin{tasks}
\task Selesaikan $(1-x^2)y''-xy' + p^2y = 0$ menggunakan metode deret pangkat di $x_0=0$.
\task Untuk nilai $p$ apa terdapat solusi polinomial?
\end{tasks}
\end{exercise}

\begin{exercise}
Carilah solusi polinomial untuk $(x^2+1) y''-2xy'+2y = 0$ dengan menggunakan
metode deret pangkat.
\end{exercise}

\begin{exercise}
\leavevmode
\begin{tasks}
\task Gunakan metode deret pangkat untuk menyelesaikan $(1-x)y''+y = 0$ di titik $x_0 = 0$.
\task Gunakan solusi dari bagian a) untuk mencari solusi
bagi $xy''+y=0$ di sekitar titik $x_0=1$.
\end{tasks}
\end{exercise}

\setcounter{exercise}{100}

\begin{exercise}
Gunakan metode deret pangkat untuk menyelesaikan $y'' + 2 x^3 y = 0$ di titik $x_0 =
0$.
\end{exercise}
\exsol{%
%\begin{equation*}
%\begin{split}
%0 = y''+2 x^3 y &= 
%\Biggl( \sum_{k=2}^\infty k\,(k-1) \, a_k x^{k-2}  \Biggr)
%+
%2 x^3
%\Biggl( \sum_{k=0}^\infty a_k x^k \Biggr)
%\\
%&=
%\Biggl( \sum_{k=2}^\infty k\,(k-1) \, a_k x^{k-2}  \Biggr)
%+
%\Biggl( \sum_{k=0}^\infty 2 a_k x^{k+3} \Biggr) .
%\\
%&=
%\Biggl( \sum_{k=0}^\infty (k+2)\,(k+1) \, a_{k+2} x^k  \Biggr)
%+
%\Biggl( \sum_{k=3}^\infty 2 a_{k-3} x^k \Biggr) .
%\\
%&=
%2 a_2 +
%6 a_3 x +
%12 a_4 x^2 +
%\Biggl( \sum_{k=3}^\infty (k+2)\,(k+1) \, a_{k+2} x^k  \Biggr)
%+
%\Biggl( \sum_{k=3}^\infty 2 a_{k-3} x^k \Biggr) .
%\end{split}
%\end{equation*}
$a_2 = 0$, $a_3 = 0$, $a_4 = 0$, relasi rekurensi (untuk $k \geq 5$): $a_k
= \frac{- 2 a_{k-5}}{k(k-1)}$,
sehingga:\\
$y(x) = a_0 + a_1 x -\frac{a_0}{10} x^5 - \frac{a_1}{15} x^6
+ \frac{a_0}{450} x^{10} + \frac{a_1}{825} x^{11}
- \frac{a_0}{47250} x^{15} - \frac{a_1}{99000} x^{16}
+
\cdots$
}

\pagebreak[2]
\begin{exercise}[menantang]
Metode deret pangkat juga berlaku untuk persamaan tak homogen.
\begin{tasks}
\task Gunakan metode deret pangkat untuk menyelesaikan $y'' - x y = \frac{1}{1-x}$
di titik $x_0 = 0$. Petunjuk: Ingat deret geometri.
\task Sekarang selesaikan untuk syarat awal $y(0)=0$, $y'(0)=0$.
\end{tasks}
\end{exercise}
\exsol{%
%\begin{equation*}
%\begin{split}
%\frac{1}{1-x} =
%\sum_{k=0}^\infty x^k
%=
%y''-x y &= 
%\Biggl( \sum_{k=2}^\infty k\,(k-1) \, a_k x^{k-2}  \Biggr)
%-
%x
%\Biggl( \sum_{k=0}^\infty a_k x^k \Biggr)
%\\
%&=
%\Biggl( \sum_{k=2}^\infty k\,(k-1) \, a_k x^{k-2}  \Biggr)
%-
%\Biggl( \sum_{k=0}^\infty a_k x^{k+1} \Biggr) .
%\\
%&=
%\Biggl( \sum_{k=0}^\infty (k+2)\,(k+1) \, a_{k+2} x^k  \Biggr)
%-
%\Biggl( \sum_{k=1}^\infty a_{k-1} x^k \Biggr) .
%\\
%&=
%2 a_2 + 
%\Biggl( \sum_{k=1}^\infty (k+2)\,(k+1) \, a_{k+2} x^k  \Biggr)
%-
%\Biggl( \sum_{k=1}^\infty a_{k-1} x^k \Biggr) .
%\end{split}
%\end{equation*}
a) $a_2 = \frac{1}{2}$, dan untuk $k \geq 3$ kita mempunyai
$a_k = \frac{a_{k-3} + 1}{k(k-1)}$, sehingga \\
$y(x) = a_0 + a_1 x + \frac{1}{2} x^2
+ \frac{a_0 + 1}{6} x^3
+ \frac{a_1 + 1}{12} x^4
+ \frac{3}{40} x^5
+ \frac{a_0 + 7}{180} x^6
+ \frac{a_1 + 13}{504} x^7
+ \frac{43}{2240} x^8
+ \frac{a_0 + 187}{12960} x^9
+ \frac{a_1 + 517}{45360} x^{10} +
\cdots$
\\
b)
$y(x) = \frac{1}{2} x^2
+ \frac{1}{6} x^3
+ \frac{1}{12} x^4
+ \frac{3}{40} x^5
+ \frac{7}{180} x^6
+ \frac{13}{504} x^7
+ \frac{43}{2240} x^8
+ \frac{187}{12960} x^9
+ \frac{517}{45360} x^{10} +
\cdots$
}

\begin{exercise}
Cobalah menyelesaikan $x^2 y'' - y = 0$ di $x_0 = 0$ menggunakan metode deret pangkat
dalam bagian ini ($x_0$ adalah titik singular).
Dapatkah Anda menemukan setidaknya satu solusi? Dapatkah Anda menemukan lebih dari satu solusi?
\end{exercise}
\exsol{%
%\begin{equation*}
%\begin{split}
%0 = x^2 y''-y &= 
%x^2 \Biggl( \sum_{k=2}^\infty k\,(k-1) \, a_k x^{k-2}  \Biggr)
%-
%\Biggl( \sum_{k=0}^\infty a_k x^k \Biggr)
%\\
%&=
%\Biggl( \sum_{k=2}^\infty k\,(k-1) \, a_k x^k  \Biggr)
%-
%a_0
%-
%a_1 x
%-
%\Biggl( \sum_{k=2}^\infty a_k x^k \Biggr) .
%\end{split}
%\end{equation*}
%sehingga $a_0 = 0$, $a_1 = 0$, $k(k-1) a_k = a_k$
Dengan menerapkan metode bagian ini secara langsung, kita memperoleh $a_k = 0$ untuk
semua $k$, sehingga $y(x) = 0$ adalah satu-satunya solusi yang kita temukan.
}

%%%%%%%%%%%%%%%%%%%%%%%%%%%%%%%%%%%%%%%%%%%%%%%%%%%%%%%%%%%%%%%%%%%%%%%%%%%%%%

\sectionnewpage
\section{Titik singular dan metode Frobenius}
\label{frobenius:section}

%mbxINTROSUBSECTION

\sectionnotes{1,5 atau 2 kuliah\EPref{, \S8.3 dan \S8.4 dalam \cite{EP}}\BDref{,
\S5.4--\S5.7 dalam \cite{BD}}}

Perilaku PDB di titik singular dapat rumit.
Untuk titik singular tertentu,
kita dapat mencari solusi setidaknya pada satu sisi titik singular
dengan menggunakan modifikasi
deret pangkat.
Mari kita tinjau beberapa contoh sebelum memberikan metode umum.
%Kita mungkin beruntung dan memperoleh solusi deret pangkat
%menggunakan metode bagian sebelumnya, tetapi secara umum kita
%harus mencoba hal-hal lain.

\subsection{Contoh}

\begin{example}
Tinjau persamaan orde pertama sederhana 
\begin{equation*}
2 x y' - y = 0 .
\end{equation*}
Perhatikan bahwa $x=0$ adalah titik singular.
Dengan menetapkan $x=0$ dalam persamaan, kita memperoleh
bahwa setiap solusi yang didefinisikan di dekat nol memenuhi
$y(0)=0$, tetapi keadaannya bahkan lebih buruk.
Jika kita mencoba mensubstitusikan
\begin{equation*}
y = \sum_{k=0}^\infty a_k x^k ,
\end{equation*}
kita memperoleh
\begin{equation*}
\begin{split}
0 = 2 xy'-y &= 
2x \, \Biggl( \sum_{k=1}^\infty k a_k x^{k-1}  \Biggr)
-
\Biggl( \sum_{k=0}^\infty a_k x^k \Biggr)
\\
& =
-a_0 + 
\sum_{k=1}^\infty (2 k a_k - a_k) \, x^{k} .
\end{split}
\end{equation*}
Pertama, $a_0 = 0$. Selanjutnya, satu-satunya cara menyelesaikan
$0 = 2 k a_k - a_k = (2k-1) \, a_k$
untuk $k = 1,2,3,\dots$ adalah dengan $a_k = 0$ untuk semua $k$.
Oleh karena itu, dengan cara ini kita hanya memperoleh solusi trivial $y=0$. Kita memerlukan
solusi tak nol untuk memperoleh solusi umum persamaan tersebut.

Mari kita coba $y=x^r$
untuk suatu bilangan real $r$. 
Akibatnya, solusi kita---jika kita dapat
menemukannya---mungkin hanya bermakna untuk $x$ positif.
Kemudian $y' = r x^{r-1}$. Jadi
\begin{equation*}
0 = 2 x y' - y = 2 x r x^{r-1} - x^r = (2r-1) x^r .
\end{equation*}
Dengan demikian, $r= \nicefrac{1}{2}$ dan karena itu $y = x^{1/2}$.
Karena persamaan tersebut linear,
solusi umum untuk $x$ positif adalah
\begin{equation*}
y = C x^{1/2} .
\end{equation*}
Jika $C \neq 0$, maka
turunan solusi \myquote{meledak} di $x=0$ (titik
singular). Hanya ada satu solusi yang terdiferensialkan
di $x=0$, yaitu solusi trivial $y=0$.
\end{example}

Tentu saja, tidak setiap masalah dengan titik singular mempunyai solusi berbentuk $y=x^r$.
Namun, mungkin kita dapat menggabungkan kedua metode tersebut. Yang akan kita lakukan adalah
mencoba solusi berbentuk
\begin{equation*}
y = x^r f(x)
\avoidbreak
\end{equation*}
untuk $x$ positif,
di mana $f(x)$ adalah fungsi analitik (deret pangkat).

\begin{example}
Tinjau persamaan
\begin{equation*}
4 x^2 y'' - 4 x^2 y' + (1-2x)y = 0,
\end{equation*}
dan perhatikan kembali bahwa $x=0$ adalah titik singular.

Mari kita coba
\begin{equation*}
y = x^r \sum_{k=0}^\infty a_k x^k
= \sum_{k=0}^\infty a_k x^{k+r} ,
\end{equation*}
di mana $r$ adalah bilangan real, tidak harus bilangan bulat.
Sekali lagi, jika solusi semacam itu ada, solusi tersebut mungkin hanya ada untuk $x$ positif.
Pertama-tama kita cari turunannya
\begin{align*}
y' & = \sum_{k=0}^\infty (k+r)\, a_k x^{k+r-1} , \\
y'' & = \sum_{k=0}^\infty (k+r)\,(k+r-1)\, a_k x^{k+r-2} .
\end{align*}
Kita substitusikan semuanya ke dalam persamaan:
\begin{equation*}
\begin{split}
0 & = 4x^2y''-4x^2y'+(1-2x)y
\\
&= 
4x^2 \, \Biggl( \sum_{k=0}^\infty (k+r)\,(k+r-1) \, a_k x^{k+r-2}  \Biggr)
%mbxlatex \\
%mbxlatex & \phantom{mmm}
-
4x^2 \, \Biggl( \sum_{k=0}^\infty (k+r) \, a_k x^{k+r-1}  \Biggr)
+
(1-2x)
\Biggl( \sum_{k=0}^\infty a_k x^{k+r} \Biggr)
\\
&=
\Biggl( \sum_{k=0}^\infty 4 (k+r)\,(k+r-1) \, a_k x^{k+r}  \Biggr)
\\
& \phantom{mmm}
-
\Biggl( \sum_{k=0}^\infty 4 (k+r) \, a_k x^{k+r+1}  \Biggr)
+
\Biggl( \sum_{k=0}^\infty a_k x^{k+r} \Biggr)
-
\Biggl( \sum_{k=0}^\infty 2a_k x^{k+r+1} \Biggr)
\\
&=
\Biggl( \sum_{k=0}^\infty 4 (k+r)\,(k+r-1) \, a_k x^{k+r}  \Biggr)
\\
& \phantom{mmm}
-
\Biggl( \sum_{k=1}^\infty 4 (k+r-1) \, a_{k-1} x^{k+r}  \Biggr)
+
\Biggl( \sum_{k=0}^\infty a_k x^{k+r} \Biggr)
-
\Biggl( \sum_{k=1}^\infty 2a_{k-1} x^{k+r} \Biggr)
\\
&=
4r(r-1) \, a_0 x^r  + a_0 x^r
%mbxlatex \\
%mbxlatex & \phantom{mmm}
+ 
\sum_{k=1}^\infty
\Bigl( 4 (k+r)\,(k+r-1) \, a_k
-
4 (k+r-1) \, a_{k-1}
+
a_k
-
2a_{k-1} \Bigr) \, x^{k+r} 
\\
&=
\bigl( 4r(r-1) + 1 \bigr) \, a_0 x^r
%mbxlatex \\
%mbxlatex & \phantom{mmm}
+ 
\sum_{k=1}^\infty
\Bigl( \bigl( 4 (k+r)\,(k+r-1) + 1 \bigr) \, a_k
-
\bigl( 4 (k+r-1) + 2 \bigr) \, a_{k-1} \Bigr) \, x^{k+r} .
\end{split}
\end{equation*}
Pertama, agar terdapat solusi kita harus mempunyai
$\bigl( 4r(r-1) + 1 \bigr) \, a_0 = 0$. Dengan mengandaikan $a_0 \neq 0$,
\begin{equation*}
4r(r-1) + 1 = 0 .
\end{equation*}
Persamaan ini disebut \emph{\myindex{persamaan indisial}}.
Persamaan indisial khusus ini
mempunyai akar ganda di $r = \nicefrac{1}{2}$.

Baiklah, sekarang kita mengetahui nilai yang harus dimiliki $r$. Pengetahuan itu kita peroleh hanya dengan melihat
koefisien $x^r$. Semua
koefisien $x^{k+r}$ lainnya juga harus bernilai nol, sehingga
\begin{equation*}
\bigl( 4 (k+r)\,(k+r-1) + 1 \bigr) \, a_k
-
\bigl( 4 (k+r-1) + 2 \bigr) \, a_{k-1} = 0 .
\end{equation*}
Jika kita mensubstitusikan $r=\nicefrac{1}{2}$ dan menyelesaikan untuk $a_k$, kita memperoleh
\begin{equation*}
a_k
=
\frac{4 (k+\nicefrac{1}{2}-1) + 2}{4 (k+\nicefrac{1}{2})\,(k+\nicefrac{1}{2}-1) + 1} \, a_{k-1}
=
\frac{1}{k} \, a_{k-1} .
\end{equation*}
Mari kita tetapkan $a_0 = 1$. Maka
\begin{equation*}
a_1 = \frac{1}{1} a_0 = 1 ,
\quad
a_2 = \frac{1}{2} a_1 = \frac{1}{2} ,
\quad
a_3 = \frac{1}{3} a_2 = \frac{1}{3 \cdot 2} ,
\quad
a_4 = \frac{1}{4} a_3 = \frac{1}{4 \cdot 3 \cdot 2} ,
\quad \cdots
\end{equation*}
Dengan mengekstrapolasi, kita melihat bahwa
\begin{equation*}
a_k = \frac{1}{k(k-1)(k-2) \cdots 3 \cdot 2} = \frac{1}{k!} .
\end{equation*}
Dengan kata lain,
\begin{equation*}
y = 
\sum_{k=0}^\infty a_k x^{k+r}
=
\sum_{k=0}^\infty \frac{1}{k!} x^{k+1/2}
=
x^{1/2}
\sum_{k=0}^\infty \frac{1}{k!} x^{k}
=
x^{1/2}
e^x .
\end{equation*}
Kita beruntung! Secara umum, kita tidak akan dapat menuliskan deret tersebut dalam
bentuk fungsi elementer.

Kita mempunyai satu solusi; sebutlah $y_1 = x^{1/2} e^x$.
Namun, bagaimana dengan solusi kedua? Jika
kita menginginkan solusi umum, kita memerlukan dua solusi bebas linear.
Memilih $a_0$ sebagai konstanta lain hanya menghasilkan kelipatan
konstan dari $y_1$, dan tidak ada $r$ lain yang dapat kita coba; kita hanya
mempunyai satu solusi untuk persamaan indisial. Nah, terdapat pangkat-pangkat $x$
di sana-sini dan kita mengambil turunan, jadi mungkin logaritma (
antiturunan dari $x^{-1}$) juga muncul. Ternyata kita perlu
mencoba solusi lain berbentuk
\begin{equation*}
y_2 = \sum_{k=0}^\infty b_k x^{k+r} + (\ln x) y_1 ,
\end{equation*}
yang dalam kasus kita adalah
\begin{equation*}
y_2 = \sum_{k=0}^\infty b_k x^{k+1/2} + (\ln x) x^{1/2} e^x .
\end{equation*}
Sekarang kita mendiferensialkan persamaan ini, mensubstitusikannya ke dalam persamaan
diferensial, lalu menyelesaikan untuk $b_k$. Perhitungan panjang pun menyusul dan kita
memperoleh suatu relasi rekurensi untuk $b_k$. Pembaca
dapat (dan sebaiknya) mencoba ini untuk memperoleh, misalnya, tiga suku pertama
\begin{equation*}
b_1 = b_0 -1 , \qquad b_2 = \frac{2b_1-1}{4} , \qquad b_3 =
\frac{6b_2-1}{18} , \qquad \ldots
\end{equation*}
Kemudian kita menetapkan $b_0$ dan memperoleh solusi $y_2$. Selanjutnya
kita menulis solusi umum sebagai $y = A y_1 + B y_2$.
\end{example}

\subsection{Metode Frobenius}

Sebelum memberikan metode umum, mari kita perjelas kapan metode ini dapat diterapkan.
Misalkan
\begin{equation*}
P(x) y'' + Q(x) y' + R(x) y = 0
\end{equation*}
adalah suatu PDB dengan polinomial
$P(x)$, $Q(x)$, dan $R(x)$.
Seperti sebelumnya, jika $P(x_0) = 0$, maka $x_0$ adalah suatu
titik singular.
Jika kita membagi dengan $P(x)$ untuk menempatkan persamaan dalam bentuk standar
$y'' + \frac{Q(x)}{P(x)} y' + \frac{R(x)}{P(x)} y = 0$,
mungkin singularitas yang muncul tidak terlalu buruk.
Lebih tepatnya, jika kedua limit
\begin{equation*}
\lim_{x \to x_0} ~ (x-x_0) \frac{Q(x)}{P(x)}
\qquad \text{dan} \qquad
\lim_{x \to x_0} ~ (x-x_0)^2 \frac{R(x)}{P(x)}
\end{equation*}
ada dan berhingga,
maka kita mengatakan bahwa $x_0$ adalah
suatu \emph{\myindex{titik singular reguler}}.

\begin{example}
Sering kali, dan untuk sisa bagian ini, $x_0 = 0$. Tinjau
\begin{equation*}
x^2y'' + x(1+x)y' + (\pi+x^2)y = 0 .
\end{equation*}
Tuliskan
\begin{align*}
& \lim_{x \to 0} ~x \frac{Q(x)}{P(x)} = 
\lim_{x \to 0} ~x \frac{x(1+x)}{x^2} = \lim_{x \to 0} ~(1+x) = 1 ,
\\
& \lim_{x \to 0} ~x^2 \frac{R(x)}{P(x)} = 
\lim_{x \to 0} ~x^2 \frac{(\pi+x^2)}{x^2} = \lim_{x \to 0} ~(\pi+x^2) = \pi .
\end{align*}
Jadi $x = 0$ adalah titik singular reguler.

Di sisi lain, jika kita membuat sedikit perubahan
\begin{equation*}
x^2y'' + (1+x)y' + (\pi+x^2)y = 0 ,
\end{equation*}
maka
\begin{equation*}
\lim_{x \to 0} ~x \frac{Q(x)}{P(x)} = 
\lim_{x \to 0} ~x \frac{(1+x)}{x^2} = \lim_{x \to 0} ~\frac{1+x}{x} =
\text{tidak ada}.
\end{equation*}
Dengan kata lain, limit tersebut tidak ada.
Titik $0$ singular, tetapi bukan titik singular reguler.
\end{example}

Sekarang kita membahas \emph{\myindex{Metode Frobenius}} umum\index{metode Frobenius}%
\footnote{Dinamai menurut matematikawan Jerman
\href{http://en.wikipedia.org/wiki/Ferdinand_Georg_Frobenius}{Ferdinand
Georg Frobenius} (1849--1917).}.
Demi kesederhanaan, kita hanya membahas metode tersebut di titik $x=0$.
Jika $x_0 \neq 0$, maka dalam solusi kita mengganti setiap $x$ dengan $(x-x_0)$.
Gagasan
utamanya adalah teorema berikut.

\begin{theorem}[Metode Frobenius]
Andaikan bahwa 
\begin{equation} \label{eq:frobeniusmethod}
P(x) y'' + Q(x) y' + R(x) y = 0
\end{equation}
mempunyai titik singular reguler di $x=0$, maka terdapat setidaknya
satu solusi berbentuk
\begin{equation*}
y = x^r \sum_{k=0}^\infty a_k x^k .
\end{equation*}
Solusi berbentuk ini disebut
\emph{\myindex{solusi tipe Frobenius}}.
\end{theorem}

\pagebreak[2]
Metode tersebut biasanya dijabarkan sebagai berikut:

\begin{enumerate}[(i)]
\item
Kita mencari solusi tipe Frobenius berbentuk
\begin{equation*}
y = \sum_{k=0}^\infty a_k x^{k+r} .
\end{equation*}
Kita substitusikan $y$ ini ke ruas kiri
persamaan \eqref{eq:frobeniusmethod}, mengumpulkan
suku-sukunya, lalu menuliskan semuanya sebagai satu deret.
\item
Deret yang diperoleh harus bernilai nol. Dengan menetapkan koefisien pertama
(yang berpangkat paling kecil) dalam deret
sama dengan nol, kita memperoleh
\emph{persamaan indisial}, yang merupakan polinomial kuadrat dalam $r$.
\item
Jika persamaan indisial mempunyai dua akar real $r_1$ dan $r_2$
sedemikian sehingga $r_1 - r_2$ bukan bilangan bulat, maka kita mempunyai dua solusi tipe
Frobenius yang bebas linear. Dengan menggunakan akar pertama, kita substitusikan
\begin{equation*}
y_1 = x^{r_1} \sum_{k=0}^\infty a_k x^{k} ,
\end{equation*}
dan kita menyelesaikan untuk semua $a_k$ guna memperoleh solusi pertama. Kemudian,
dengan menggunakan akar kedua,
kita substitusikan
\begin{equation*}
y_2 = x^{r_2} \sum_{k=0}^\infty b_k x^{k} ,
\end{equation*}
dan menyelesaikan untuk semua $b_k$ guna memperoleh solusi kedua.
\item
Jika persamaan indisial mempunyai akar ganda $r$, maka kita menemukan 
satu solusi
\begin{equation*}
y_1 = x^{r} \sum_{k=0}^\infty a_k x^{k} ,
\end{equation*}
kemudian kita memperoleh solusi baru dengan mensubstitusikan
\begin{equation*}
y_2 = x^{r} \sum_{k=0}^\infty b_k x^{k} + (\ln x) y_1 ,
\end{equation*}
ke dalam persamaan \eqref{eq:frobeniusmethod} dan menyelesaikan untuk konstanta-konstanta $b_k$.
\item
Jika persamaan indisial mempunyai dua akar real sedemikian sehingga $r_1-r_2$
adalah bilangan bulat, maka, dengan mengandaikan $r_1$ adalah akar yang lebih besar, satu solusinya adalah
\begin{equation*}
y_1 = x^{r_1} \sum_{k=0}^\infty a_k x^{k} ,
\end{equation*}
dan solusi bebas linear kedua berbentuk
\begin{equation*}
y_2 = x^{r_2} \sum_{k=0}^\infty b_k x^{k} + C (\ln x) y_1 ,
\end{equation*}
di mana kita mensubstitusikan $y_2$ ke dalam \eqref{eq:frobeniusmethod} dan menyelesaikan untuk
konstanta-konstanta $b_k$ dan $C$.
\item
Terakhir, jika persamaan indisial mempunyai akar kompleks, maka penyelesaian
untuk $a_k$ dalam solusi
\begin{equation*}
y = x^{r_1} \sum_{k=0}^\infty a_k x^{k}
\end{equation*}
menghasilkan fungsi bernilai kompleks---$a_k$ adalah bilangan
kompleks. Kita memperoleh dua solusi bebas
linear kita%
\footnote{Lihat 
Joseph L.\ Neuringer,
\emph{The Frobenius method for complex roots of the indicial equation},
International Journal of Mathematical Education in Science and Technology,
Volume 9, Issue 1, 1978, 71--77.}
dengan mengambil bagian real dan imajiner dari $y$.
\end{enumerate}

Gagasan utamanya adalah mencari setidaknya satu solusi tipe Frobenius. Jika
kita beruntung dan menemukan dua, pekerjaan selesai.
Jika kita hanya memperoleh satu, kita menggunakan gagasan di atas atau bahkan metode lain
seperti reduksi orde (lihat \sectionref{solinear:section}) untuk
memperoleh solusi kedua.

\subsection{Fungsi Bessel} \label{bessel:subsection}

Suatu kelas fungsi penting yang umum muncul dalam fisika adalah
\emph{fungsi Bessel}%
\footnote{Dinamai menurut
astronom dan matematikawan Jerman
\href{http://en.wikipedia.org/wiki/Friedrich_Bessel}{Friedrich Wilhelm
Bessel} (1784--1846).}.
Sebagai contoh, fungsi-fungsi ini muncul ketika menyelesaikan
persamaan gelombang dalam dua dan tiga dimensi. Pertama, tinjau
\emph{\myindex{persamaan Bessel}} berorde $p$:
\begin{equation*}
x^2 y'' + xy' + \left(x^2 - p^2\right)y = 0 .
\end{equation*}
Kita mengizinkan $p$ sebagai sembarang bilangan, bukan hanya bilangan bulat, meskipun bilangan bulat
dan kelipatan $\nicefrac{1}{2}$ paling penting dalam penerapan.

Ketika kita mensubstitusikan
\begin{equation*}
y = \sum_{k=0}^\infty a_k x^{k+r}
\end{equation*}
ke dalam persamaan Bessel berorde $p$, kita memperoleh persamaan indisial
\begin{equation*}
r(r-1)+r-p^2 = (r-p)(r+p) = 0 .
\end{equation*}
Kecuali jika $p=0$, kita memperoleh dua akar, $r_1 = p$ dan $r_2 = -p$.
Jika $p$ bukan bilangan bulat, maka dengan mengikuti metode Frobenius dan
menetapkan $a_0 = 1$, kita memperoleh
solusi-solusi bebas linear berbentuk
\begin{align*}
& y_1 = x^p \sum_{k=0}^\infty
\frac{{(-1)}^k x^{2k}}{2^{2k} k! (k+p)(k-1+p)\cdots (2+p)(1+p)} ,
\\
& y_2 = x^{-p} \sum_{k=0}^\infty
\frac{{(-1)}^k x^{2k}}{2^{2k} k! (k-p)(k-1-p)\cdots (2-p)(1-p)} .
\end{align*}

\begin{exercise}
\leavevmode
\begin{tasks}
\task
Verifikasi bahwa persamaan indisial dari persamaan Bessel berorde $p$ adalah
$(r-p)(r+p)=0$.
\task
Andaikan $p$ bukan bilangan bulat. Lakukan perhitungannya
untuk mencari solusi $y_1$ dan $y_2$ di atas.
\end{tasks}
\end{exercise}

Fungsi Bessel adalah kelipatan konstan yang sesuai dari $y_1$ dan $y_2$.
Pertama, kita mendefinisikan \emph{fungsi gamma}
\begin{equation*}
\Gamma(x) = \int_0^\infty t^{x-1} e^{-t} \, dt .
\end{equation*}
Perhatikan bahwa $\Gamma(1) = 1$.
Fungsi gamma juga mempunyai sifat yang sangat baik
\begin{equation*}
\Gamma(x+1) = x \Gamma(x) .
\end{equation*}
Dari sifat ini, diperoleh bahwa $\Gamma(n) = (n-1)!$ ketika $n$ adalah
bilangan bulat positif. Jadi fungsi gamma adalah versi kontinu dari faktorial. Kita
menghitung:
\begin{align*}
& \Gamma(k+p+1)=(k+p)(k-1+p)\cdots (2+p)(1+p) \Gamma(1+p) ,
\\
& \Gamma(k-p+1)=(k-p)(k-1-p)\cdots (2-p)(1-p) \Gamma(1-p) .
\end{align*}

\begin{exercise}
Verifikasi identitas-identitas di atas dengan menggunakan 
$\Gamma(x+1) = x \Gamma(x)$.
\end{exercise}

Dengan menggunakan $\Gamma(x+1) = x \Gamma(x)$,
kita mendefinisikan fungsi gamma untuk 
semua $x$ kecuali bilangan bulat tak positif. 
Kebalikan $\frac{1}{\Gamma(x)}$
dapat didefinisikan untuk semua $x$ dengan menetapkan $\frac{1}{\Gamma(x)}$ bernilai nol
untuk bilangan bulat $x \leq 0$.

Kita mendefinisikan \emph{fungsi Bessel jenis pertama}%
\index{fungsi Bessel jenis pertama} berorde
$p$ dan $-p$ sebagai
\begin{align*}
& J_p(x) = \frac{1}{2^p\Gamma(1+p)} y_1
=
\sum_{k=0}^\infty
\frac{{(-1)}^k}{k! \, \Gamma(k+p+1)}
{\left(\frac{x}{2}\right)}^{2k+p} ,
\\
& J_{-p}(x) = \frac{1}{2^{-p}\Gamma(1-p)} y_2
=
\sum_{k=0}^\infty
\frac{{(-1)}^k}{k! \,\Gamma(k-p+1)}
{\left(\frac{x}{2}\right)}^{2k-p} .
\end{align*}
Karena fungsi-fungsi ini merupakan kelipatan konstan dari solusi yang kita temukan di atas, keduanya
adalah solusi persamaan Bessel berorde $p$. Konstanta-konstanta tersebut dipilih
demi kemudahan.
Karena kebalikan dari $\Gamma$ didefinisikan untuk semua $x$, rumus-rumus
di ruas kanan berlaku bahkan untuk $p$ bilangan bulat.
Ketika $n$ adalah bilangan bulat tak negatif,
\begin{equation*}
J_n(x) =
\sum_{k=0}^\infty
\frac{{(-1)}^k}{k! \,(k+n)!}
{\left(\frac{x}{2}\right)}^{2k+n} .
\end{equation*}

Ketika $p$ bukan bilangan bulat, $J_p$
dan $J_{-p}$ bebas linear.
Untuk orde bilangan bulat, dapat diperiksa bahwa
\begin{equation*}
J_n(x) = {(-1)}^nJ_{-n}(x) .
\end{equation*}
Jadi $J_{-n}$ bukan solusi bebas linear
kedua. Solusi lainnya adalah apa yang disebut
\emph{\myindex{fungsi Bessel jenis kedua}}. Fungsi-fungsi ini hanya
bermakna untuk orde bilangan bulat $n$ dan
didefinisikan sebagai limit kombinasi linear $J_p(x)$ dan $J_{-p}(x)$, ketika
$p$ mendekati $n$ dengan
cara berikut:
\begin{equation*}
Y_n(x) = \lim_{p\to n} \frac{\cos(p \pi) J_p(x) - J_{-p}(x)}{\sin(p \pi)} .
\end{equation*}
Setiap kombinasi linear $J_p(x)$ dan $J_{-p}(x)$ adalah solusi
persamaan Bessel berorde $p$. Kemudian, ketika kita mengambil limit saat $p$
menuju $n$, kita melihat bahwa $Y_n(x)$ adalah solusi persamaan Bessel
berorde $n$. Ternyata $Y_n(x)$ dan $J_n(x)$ juga bebas
linear. Oleh karena itu, ketika $n$ adalah bilangan bulat, kita mempunyai
solusi umum persamaan Bessel berorde $n$:
\begin{equation*}
y = A J_n(x) + B Y_n(x) ,
\end{equation*}
untuk konstanta sembarang $A$ dan $B$. Perhatikan bahwa
$Y_n(x)$ menuju negatif tak hingga di $x=0$. Banyak paket perangkat lunak
matematika telah mendefinisikan
fungsi-fungsi $J_n(x)$ dan $Y_n(x)$ ini, sehingga keduanya dapat digunakan
seperti, misalnya, $\sin(x)$ dan $\cos(x)$. Bahkan, fungsi Bessel
mempunyai beberapa sifat serupa. Sebagai contoh, $-J_1(x)$ adalah turunan
dari $J_0(x)$, dan secara umum turunan $J_n(x)$ dapat dituliskan sebagai kombinasi
linear dari $J_{n-1}(x)$ dan $J_{n+1}(x)$. Selain itu, fungsi-fungsi ini
berosilasi, meskipun tidak periodik. Lihat
\figurevref{bessel:graphsfig} untuk grafik fungsi Bessel.
\begin{myfig}
\capstart
%berkas asli bessel-first bessel-second
\diffyincludegraphics{width=6.24in}{width=9in}{bessel-first-second}{%
Dua grafik ditampilkan. Pada grafik sebelah kiri, dua fungsi yang tampak
berosilasi digambar untuk x dari 0 hingga 10; salah satunya bermula di
0 dengan kemiringan positif,
dan yang lain bermula di 1 dengan kemiringan nol. Amplitudonya
tampak menurun ketika x semakin besar. Pada grafik sebelah kanan,
ditampilkan dua grafik serupa, tetapi sekarang keduanya tampak bermula
dari negatif tak hingga dan mula-mula naik sepanjang sumbu y.}
\caption{Grafik $J_0(x)$ dan $J_1(x)$ pada grafik pertama
serta $Y_0(x)$ dan $Y_1(x)$ pada grafik kedua.\label{bessel:graphsfig}}
\end{myfig}

\begin{example}
Persamaan lain kadang-kadang dapat diselesaikan dalam bentuk fungsi Bessel.
Sebagai contoh, diberikan konstanta positif $\lambda$,
\begin{equation*}
x y'' + y' + \lambda^2 x y = 0 ,
\end{equation*}
dapat diubah menjadi 
$x^2 y'' + x y' + \lambda^2 x^2 y = 0$. Dengan mengganti variabel
$t = \lambda x$, melalui aturan rantai kita memperoleh persamaan dalam $y$ dan $t$:
\begin{equation*}
t^2 y'' + t y' + t^2 y = 0 ,
\end{equation*}
yang kita kenali sebagai persamaan Bessel berorde $0$. Oleh karena itu,
solusi umumnya adalah $y(t) = A J_0(t) + B Y_0(t)$, atau dalam $x$:
\begin{equation*}
y = A J_0(\lambda x) + B Y_0(\lambda x) .
\end{equation*}
Persamaan ini muncul, misalnya, ketika mencari modus dasar
getaran suatu drum melingkar, tetapi kita telah menyimpang dari pembahasan.
\end{example}

\subsection{Latihan}

\begin{exercise}
Carilah suatu solusi khusus (tipe Frobenius) dari $x^2 y'' + x y' + (1+x) y = 0$.
\end{exercise}

\begin{exercise}
Carilah suatu solusi khusus (tipe Frobenius) dari $x y'' - y = 0$.
\end{exercise}

\begin{exercise}
Carilah suatu solusi khusus (tipe Frobenius) dari $y'' +\frac{1}{x}y' - xy = 0$.
\end{exercise}

\begin{exercise}
Carilah solusi umum dari $2 x y'' + y' - x^2 y = 0$.
\end{exercise}

\begin{exercise}
Carilah solusi umum dari $x^2 y'' - x y' -y = 0$.
\end{exercise}

\begin{exercise}
Dalam persamaan-persamaan berikut,
klasifikasikan titik $x=0$ sebagai \emph{biasa}, \emph{singular reguler}, atau
\emph{singular tetapi tidak singular reguler}.
\begin{tasks}(2)
\task $x^2(1+x^2)y''+xy=0$
\task $x^2y''+y'+y=0$
\task $xy''+x^3y'+y=0$
\task $xy''+xy'-e^xy=0$
\task $x^2y''+x^2y'+x^2y=0$
\end{tasks}
\end{exercise}

\setcounter{exercise}{100}

\begin{exercise}
Dalam persamaan-persamaan berikut,
klasifikasikan titik $x=0$ sebagai \emph{biasa}, \emph{singular reguler}, atau
\emph{singular tetapi tidak singular reguler}.
\begin{tasks}(2)
\task $y''+y=0$
\task $x^3y''+(1+x)y=0$
\task $xy''+x^5y'+y=0$
\task $\sin(x)y''-y=0$
\task $\cos(x)y''-\sin(x)y=0$
\end{tasks}
\end{exercise}
\exsol{%
a)~biasa, b)~singular tetapi tidak singular reguler, c)~singular reguler,
d)~singular reguler, e)~biasa.
}

\begin{exercise}
Carilah solusi umum dari $x^2 y'' -y = 0$.
\end{exercise}
\exsol{%
$y = A x^{\frac{1+\sqrt{5}}{2}} + B x^{\frac{1-\sqrt{5}}{2}}$
}

\begin{exercise}
Carilah suatu solusi khusus dari $x^2 y'' +(x-\nicefrac{3}{4})y = 0$.
\end{exercise}
\exsol{%
$y = x^{3/2}
\sum\limits_{k=0}^\infty
\frac{{(-1)}^{k}}{k!\,(k+2)!} x^k$ (Perhatikan bahwa demi kemudahan,
kita tidak memilih $a_0 = 1$.)}

\begin{exercise}[rumit]
Carilah solusi umum dari $x^2 y'' - x y' +y = 0$.
\end{exercise}
\exsol{%
$y = Ax + B x \ln(x)$
}
